% !PRINTED 134
\weberpage{134}

Um den zweiten Theil zu beweisen, setzen wir
\begin{align}
 f(x)&=\psi(x)(x-\alpha)^a,
\tag{8}\\
 \psi(x)&=(x-\beta)^b(x-\gamma)^c\cdots .
\tag{9}
\end{align}
Wir wählen nun eine positive Zahl $A$, die aber von $\varrho$ unabhängig sein soll, so dass, so lange $x$ im Inneren oder an der Peripherie des Gebietes $(\varrho)$ liegt,
\begin{equation}
 |\psi(x)|>A \qquad (\hbox{für } |x-\alpha|\le \varrho)
\tag{10}
\end{equation}
bleibt. Eine solche Zahl erhalten wir z. B., wenn wir $l$ gleich der Hälfte des kleinsten unter den Abständen $(\alpha,\beta),(\alpha,\gamma),\ldots$ annehmen und
\[
 A=l^{b+c+\cdots}
\]
setzen; dann ist die in (10) ausgedrückte Forderung wenigstens so lange erfüllt, als $\varrho$ kleiner als $l$ ist. Ist nur ein einziger Punkt $\alpha$ vorhanden, so ist $\psi(x)=1$ zu setzen, und für $A$ kann jeder beliebige echte Bruch gesetzt werden.

Es seien nun $\alpha_1,\alpha_2,
\ldots$ die Wurzeln von $f_1(x)$ innerhalb $(\varrho)$, und $a_1,a_2,
\ldots$ die Grade ihrer Vielfachheit, und $\beta_1,\beta_2,
\ldots$, $b_1,b_2,
\ldots$, $\gamma_1,\gamma_2,
\ldots$, $c_1,c_2,
\ldots$ sollen dieselbe Bedeutung für die Gebiete $(\varrho'),(\varrho''),\ldots$ haben; es ist dann
\begin{equation}
 n=a_1+a_2+
\cdots+b_1+b_2+\cdots+c_1+c_2+\cdots,
\tag{11}
\end{equation}
da die Gesammtzahl aller Wurzeln gleich $n$ sein muss. Die Function $f_1(x)$ lässt sich dann so darstellen:
\begin{equation}
 f_1(x)=\psi_1(x)(x-\alpha_1)^{a_1}(x-\alpha_2)^{a_2}\cdots,
\tag{12}
\end{equation}
worin
\begin{equation}
 \psi_1(x)=(x-\beta_1)^{b_1}(x-\beta_2)^{b_2}\cdots
          (x-\gamma_1)^{c_1}(x-\gamma_2)^{c_2}\cdots .
\tag{13}
\end{equation}

Nun können wir eine von $\varrho,\varrho',\varrho''$ unabhängige positive Zahl $B$ bestimmen, so dass, so lange $x$ innerhalb oder an der Grenze von $(\varrho)$ bleibt,
\begin{equation}
 |\psi_1(x)|<B \qquad (\hbox{für } |x-\alpha|\le \varrho).
\tag{14}
\end{equation}
Wir können z. B. eine Grösse $L$ wählen, die grösser ist als die doppelte Entfernung des Punktes $\alpha$ von einem der Punkte $\beta,\gamma,
\ldots$ und jedenfalls grösser als $1$ und dann
\[
 B>L^n
\]
annehmen.

Nun folgt aus (4)
% [sic] The source prints a strict triangle inequality in (15); equality is possible, so the universally valid relation is <=. The later strict bounds still suffice for the argument.
\begin{equation}
 |f(x)|<|f_1(x)|+|\varphi(x)|,
\tag{15}
\end{equation}
%CONT>>
